% TODO make hw stylesheet
\documentclass[10pt]{scrartcl}
\usepackage[a4paper, total={6in, 8in}]{geometry}
\usepackage[usenames,dvipsnames,svgnames]{xcolor}
\usepackage[shortlabels]{enumitem}
\usepackage[framemethod=TikZ]{mdframed}
\usepackage{amsmath,amssymb,amsthm}
\usepackage[colorlinks]{hyperref}
\usepackage{multicol}
\usepackage{tabularx}

\hypersetup{
	colorlinks=true,
	linkcolor=blue,
	filecolor=magenta,      
	urlcolor=magenta,
	pdftitle={MAT 322 HW}
}

\usepackage{microtype}
\usepackage{mathtools}
\usepackage[headsepline]{scrlayer-scrpage}
\usepackage{thmtools}
\usepackage[textsize=footnotesize]{todonotes}

\mdfdefinestyle{mdbluebox}{roundcorner=10pt,innerbottommargin=9pt,
	linecolor=blue,backgroundcolor=TealBlue!5,}
\declaretheoremstyle[headfont=\sffamily\bfseries\color{MidnightBlue},
mdframed={style=mdbluebox},]{thmbluebox}

\mdfdefinestyle{mdbloobox}{frametitlefont=\bfseries,innerbottommargin=8pt,
	nobreak=true,backgroundcolor=TealBlue!5,linecolor=blue,}
\declaretheoremstyle[headfont=\bfseries\color{MidnightBlue},
mdframed={style=mdbloobox},headpunct={\\[3pt]},postheadspace=0pt,]{thmbloobox}

\mdfdefinestyle{mdredbox}{frametitlefont=\bfseries,innerbottommargin=8pt,
	nobreak=true,backgroundcolor=Salmon!5,linecolor=RawSienna,}
\declaretheoremstyle[headfont=\bfseries\color{RawSienna},
mdframed={style=mdredbox},headpunct={\\[3pt]},postheadspace=0pt,]{thmredbox}

\mdfdefinestyle{mdgreenbox}{linecolor=ForestGreen,backgroundcolor=ForestGreen!5,
	linewidth=2pt,rightline=false,leftline=true,topline=false,bottomline=false,}
\declaretheoremstyle[headfont=\bfseries\sffamily\color{ForestGreen!70!black},
mdframed={style=mdgreenbox},headpunct={ --- },]{thmgreenbox}

\mdfdefinestyle{mdorangebox}{linecolor=RedOrange,backgroundcolor=RedOrange!5,
	linewidth=2pt,rightline=false,leftline=true,topline=false,bottomline=false,}
\declaretheoremstyle[headfont=\bfseries\sffamily\color{RedOrange!70!black},
mdframed={style=mdorangebox},headpunct={ --- },]{thmorangebox}

\mdfdefinestyle{mdblackbox}{linecolor=black,backgroundcolor=RedViolet!5!gray!5,
	linewidth=3pt,rightline=false,leftline=true,topline=false,bottomline=false,}
\declaretheoremstyle[mdframed={style=mdblackbox}]{thmblackbox}

\declaretheoremstyle[spaceabove=6pt,spacebelow=6pt]{basehead}

\declaretheorem[style=basehead,name=Problem]{problem}
\declaretheorem[style=thmredbox,name=Problem,numbered=no]{problem*}
\declaretheorem[style=thmbloobox,name=Setup]{setup} % for config geo
\declaretheorem[style=thmredbox,name=Example,sibling=problem]{example}
\declaretheorem[style=thmredbox,name=$\bigstar$ Problem,sibling=problem]{niceprob}
\declaretheorem[style=thmbluebox,name=Theorem,sibling=problem]{theorem}
\declaretheorem[style=thmbluebox,name=Corollary,sibling=theorem]{corollary}
\declaretheorem[style=thmbluebox,name=Proposition,sibling=theorem]{prop}
\declaretheorem[style=thmbluebox,name=Lemma,numberwithin=problem]{lemma}
\declaretheorem[style=thmbluebox,name=Theorem,numbered=no]{theorem*}
\declaretheorem[style=thmbluebox,name=Lemma,numbered=no]{lemma*}
\declaretheorem[style=thmgreenbox,name=Claim,numberwithin=problem]{claim}
\declaretheorem[style=thmgreenbox,name=Claim,numbered=no]{claim*}
\declaretheorem[style=thmblackbox,name=Remark,numberwithin=problem]{remark}
\declaretheorem[style=thmblackbox,name=Remark,numbered=no]{remark*}
\declaretheorem[style=thmgreenbox,name=Definition,sibling=theorem]{definition}
\declaretheorem[style=thmgreenbox,name=Definition,numbered=no]{definition*}
\declaretheorem[style=thmorangebox,name=Warning,sibling=theorem]{warning}
\declaretheorem[style=thmorangebox,name=Warning,numbered=no]{warning*}
\declaretheorem[style=thmorangebox,name=Note,numbered=no]{note}
\declaretheorem[style=thmblackbox,name=Exercise,sibling=problem]{exercise}
\declaretheorem[style=thmblackbox,name=Exercise,numbered=no]{exercise*}
\declaretheorem[style=thmblackbox,name=Question,sibling=problem]{question}
\declaretheorem[style=thmblackbox,name=Question,numbered=no]{question*}

\addtolength{\textheight}{3.14cm}
\providecommand{\ol}{\overline}
\providecommand{\ul}{\underline}
\providecommand{\eps}{\varepsilon}
\providecommand{\half}{\frac{1}{2}}
\providecommand{\dang}{\measuredangle} %% Directed angle
\providecommand{\CC}{\mathbb C}
\providecommand{\FF}{\mathbb F}
\providecommand{\NN}{\mathbb N}
\providecommand{\QQ}{\mathbb Q}
\providecommand{\RR}{\mathbb R}
\providecommand{\ZZ}{\mathbb Z}
\providecommand{\dg}{^\circ}
\providecommand{\ii}{\item}
\providecommand{\setdiff}{\smallsetminus}
\providecommand{\alert}[1]{{\sffamily\textbf{\textcolor{blue}{#1}}}}
\providecommand{\inv}{^{-1}}
\providecommand{\mb}[1]{\mathbf{#1}}
\providecommand{\dif}{\;d}

\reversemarginpar

\newenvironment{soln}{\begin{proof}[Solution]}{\end{proof}}
\newenvironment{walkthrough}{\noindent \textbf{Walkthrough.}}{\medskip}
\newlist{walk}{enumerate}{3}
\setlist[walk]{label=\bfseries (\alph*)}

\providecommand{\pow}{\operatorname{Pow}}
\providecommand{\vocab}[1]{{\textbf{\textcolor{ForestGreen}{#1}}}}
\providecommand{\lcm}{\operatorname{lcm}}
\providecommand{\abs}[1]{\left|#1\right|}

\usepackage{asymptote}
\begin{asydef}
	defaultpen(fontsize(10pt));
	unitsize(1cm);
	size(8cm); // set a reasonable default
	usepackage("amsmath");
	usepackage("amssymb");
	settings.tex="pdflatex";
	settings.outformat="pdf";
	// Replacement for olympiad+cse5 which is not standard
	import geometry;
	// recalibrate fill and filldraw for conics
	void filldraw(picture pic = currentpicture, conic g, pen fillpen=defaultpen, pen drawpen=defaultpen)
	{ filldraw(pic, (path) g, fillpen, drawpen); }
	void fill(picture pic = currentpicture, conic g, pen p=defaultpen)
	{ filldraw(pic, (path) g, p); }
	// some geometry
	pair foot(pair P, pair A, pair B) { return foot(triangle(A,B,P).VC); }
	pair orthocenter(pair A, pair B, pair C) { return orthocentercenter(A,B,C); }
	pair centroid(pair A, pair B, pair C) { return (A+B+C)/3; }
	// cse5 abbreviations
	path CP(pair P, pair A) { return circle(P, abs(A-P)); }
	path CR(pair P, real r) { return circle(P, r); }
	pair IP(path p, path q) { return intersectionpoints(p,q)[0]; }
	pair OP(path p, path q) { return intersectionpoints(p,q)[1]; }
	path Line(pair A, pair B, real a=0.6, real b=a) { return (a*(A-B)+A)--(b*(B-A)+B); }
	// cse5 more useful functions
	picture CC() {
		picture p=rotate(0)*currentpicture;
		currentpicture.erase();
		return p;
	}
	pair MP(Label s, pair A, pair B = plain.S, pen p = defaultpen) {
		Label L = s;
		L.s = "$"+s.s+"$";
		label(L, A, B, p);
		return A;
	}
	pair Drawing(Label s = "", pair A, pair B = plain.S, pen p = defaultpen) {
		dot(MP(s, A, B, p), p);
		return A;
	}
	path Drawing(path g, pen p = defaultpen, arrowbar ar = None) {
		draw(g, p, ar);
		return g;
	}
\end{asydef}

\usepackage{mathrsfs}

\ihead{\footnotesize MAT 322 HW07}
\ohead{\footnotesize Last compiled \today}

\title{\vspace{-2em}MAT 322 HW07}
\author{\large Aditya Pahuja}
\date{\large Due 04/08}
\begin{document}
\maketitle

\begin{problem}
	Let $U\subseteq\RR^n$ be open and let $f\colon U\to\RR$ be continuous.
	If $f$ vanishes outside a compact subset $C\subseteq U$,
	show that $\int_Uf$ and $\int_Cf$ both exist and are equal.
\end{problem}

\begin{soln}
	Let $C_1\subseteq C_2\subseteq\cdots\subseteq U$ be a sequence
	of compact subsets of $U$ whose union is $U$.
	Let $C'_i = C_i\cup C$, so that now the sequence $C'_i$
	satisfies the same properties as $C_i$, but with each $C'_i$ containing $C$.
	Now clearly
	\[\int_{C'_i}f = \int_Cf \le \max_{x\in C}f \cdot\operatorname{vol}C
	< \infty,\]
	with existence of the integrals given by being nowhere discontinuous. Then
	\[\int_Uf = \lim_{i\to\infty}\int_{C'_i}f = \int_Cf\]
	as desired.
\end{soln}

\begin{problem}
	Let $U$ be the region in $\RR^2$ bounded by the curve $x^2-xy+2y^2=1$.
	Express the integral $\int_Uxy$ as an integral
	over the unit ball in $\RR^2$ centered at $0$.
\end{problem}

\begin{soln}
	Let $a = \frac{\sqrt{2-\sqrt2}}{2}$
	and $b = -\frac{\sqrt{2+\sqrt2}}{2}$,
	and let $s = \frac{3+\sqrt2}{2}$ and $t = \frac{3-\sqrt2}{2}$.
	One can then compute that the map
	\[f(x,y) := \left(\frac{ax}{\sqrt s} - \frac{by}{\sqrt t},
	\frac{bx}{\sqrt s} + \frac{ay}{\sqrt t}\right)\]
	bijectively sends the unit ball $B$ centered at $0$ to the region bounded by
	the ellipse $x^2-xy+2y^2=1$. Therefore, by the change of variables formula,
	\begin{align*}
		\int_{f(B)} xy &= \int_{B} g\circ f\cdot |\det(\mathrm df)|\\
		&= \int_B \left(\frac{abx^2}{s} + \frac{(a^2-b^2)xy}{\sqrt{st}}
		- \frac{aby^2}{t}\right)\cdot\frac{a^2+b^2}{\sqrt{st}}\\
		&= \int_B \frac{4}{7}\cdot\left(\frac{\sqrt2}{2(3+\sqrt2)}x^2
		- \frac{2\sqrt2}{7}xy-\frac{\sqrt2}{2(3-\sqrt2)}y^2\right)
	\end{align*}
	where $g(x,y) = xy$.
\end{soln}

\begin{problem}
	\begin{enumerate}[(a)]
		\ii Express the volume of the solid in $\RR^3$ bounded below
		by the surface $z=x^2+2y^2$ and above by the plane $z=2x+6y+1$,
		as the integral of a suitable function over
		the unit ball in $\RR^2$ centered at $0$.
		\ii Compute this volume.
	\end{enumerate}
\end{problem}

\begin{soln}
	The surface and plane intersect on the locus of points satisfying
	\[x^2-2x+2y^2-6y-1=0\iff
	(x-1)^2 + 2\left(y-\frac{3}{2}\right)^2 = \frac{13}{2}.\]
	Let the region bounded by this ellipse be $E$; we want to compute
	\[\int_E 2x-x^2+6y-2y^2+1.\]
	Translate $E$ to get ellipse $E'$ centered at the origin first,
	so the integral becomes
	\[\int_{E'} \frac{13}{2} - x^2 - 2y^2.\]
	Then, take the map $(x,y)\mapsto(x\sqrt{2/13}, y\sqrt{4/13})$
	to shrink $E'$ to the unit disk, which gives
	\[\int_{B^2(1)}\left(\frac{13}{2}-\frac{13x^2}{2} - \frac{13y^2}{2}\right)
	\cdot \frac{2\sqrt2}{13}
	= \int_{\theta=0}^{2\pi}\int_{r=0}^1 r^3\sqrt2
	= \boxed{\frac{\pi\sqrt2}{2}}.\]
\end{soln}

\begin{problem}
	Let $B^n(R)$ denote the closed ball of radius $R$ in $\RR^n$.
	\begin{enumerate}[(a)]
		\ii Show that $\operatorname{vol}(B^n(R)) = \lambda_nR^n$
		for some constant $\lambda_n$.
		\ii Compute $\lambda_1$ and $\lambda_2$.
		\ii Compute $\lambda_n$ in terms of $\lambda_{n-2}$.
		\ii Obtain a formula for $\lambda_n$.
	\end{enumerate}
\end{problem}

\begin{soln}
	(a) Let $\lambda_n$ be the volume of the unit ball $B^n(1)$.
	Then, we have
	\[\int_{B^n(R)} 1
	= \int_{B^n(1)} R^n = \lambda_nR^n\]
	using the change of variables formula with the scaling-by-factor-of-$R$ map.
	
	(b) The volume of $B^1(1) = (-1,1)$ is $\boxed{\lambda_1=2}$,
	and the volume of $B^2(1)$ is the area of the unit circle
	$\boxed{\lambda_2=\pi}$.
	
	(c) We can compute the volume by integrating
	the $n-2$ dimensional balls above each point in the $x_1x_2$-plane:
	specifically, if one picks a point a distance of $r \le 1$ from $0$,
	then the cross-section of the sphere obtained by fixing
	$x_1$ and $x_2$ is $B^{n-2}(\sqrt{R^2-r^2})$. Therefore,
	parametrizing the $x_1x_2$-plane in polar coordinates,
	the volume of the $n$-dimensional unit ball is
	\[\lambda_n = \int_{\theta=0}^{2\pi} \int_{r=0}^1
	\lambda_{n-2}\sqrt{1-r^2}^{n-2}\cdot r
	= 2\pi\lambda_{n-2}\cdot\left[-\frac{1}{n}(1-r^2)^{n/2}\right]_{r=0}^1
	= \frac{2\pi\lambda_{n-2}}{n}.\]
	
	(d) If $n$ is even, we get
	\[\lambda_n = \frac{2\pi}{n}\cdot\lambda_{n-2}
	= \frac{2\pi}{n}\cdot\frac{2\pi}{n-2}
	= \cdots = \frac{(2\pi)^{n/2}}{n!!}\]
	and if $n$ is odd,
	\[\lambda_n = \frac{2\pi}{n}\cdot\lambda_{n-2}
	= \frac{2\pi}{n}\cdot\frac{2\pi}{n-2}
	= \cdots = \frac{(2\pi)^{(n-1)/2}}{n!!}\cdot2\]
	where $n!!$ is the product of all positive integers
	less than or equal to $n$ that have the same parity as $n$
	(so for example $5!! = 5\cdot3\cdot 1$ and $4!! = 4\cdot2$).
\end{soln}


















\end{document}